\chapter{Discussion and Conclusions}

% \section{Summary of experimental results}

% % TODO: write this section.

\section{Strengths and limitations of Compatible Substitutions Optimization}

The main strength of CoSO is that it performs representation-aware search directly in the original genotype space. Unlike latent-space methods, it does not require training a generative model, selecting checkpoints, or decoding continuous vectors back into genotype strings. Candidate changes are expressed as explicit fragment substitutions in \fone{}, and every accepted modification can be inspected as a replacement of one genotype fragment by another. This makes the method comparatively transparent: the role of the archive, compatibility thresholds, accepted substitutions, and local convergence events can be analysed directly from the run logs.

A second strength is that CoSO can reuse information gathered during search. Fragments extracted from locally converged genotypes are stored in the archive and can later become candidate replacements in different genotype contexts. The bucket-aligned results indicate that this mechanism was most useful when the starting genotype pool left substantial room for improvement. In the weaker and medium-fitness buckets, the best CoSO configurations produced clear improvements beyond the best introduced genotype. This suggests that explicit compatible substitutions can provide useful search moves without relying on a learned continuous representation.

The method also offers interpretable control through its compatibility thresholds. The spatial one determines how similar the geometric displacement of two fragments must be, while the directional threshold limits differences between their terminal directions. The experiments showed that this control is important: overly weak or poorly matched spatial compatibility produced much smaller improvements, whereas the stronger configurations benefited from more effective fragment reuse. CoSO therefore exposes a meaningful trade-off between conservative fragment matching and broader archive retrieval.

The most important limitation is that CoSO depends on hand-designed representation features. In this thesis, compatibility is defined using spatial and directional descriptors derived from Framsticks \fone{} fragments. These descriptors are meaningful for the studied representation, but they are not automatically transferable to arbitrary genotype languages or tasks. Applying CoSO to a different domain would require designing new fragment extraction rules, feature descriptors, and compatibility thresholds.

A second limitation follows from the archive-based multi-start structure of the algorithm. A CoSO run is not a single-source optimization trajectory. It may process several genotypes from the input pool while retaining an archive accumulated from previous locally converged solutions. This makes CoSO suitable for bucket-level and configuration-level analysis, but it prevents ordinary paired per-source comparisons with methods such as LSO or a standard evolutionary baseline. The best-so-far histories should therefore be interpreted as progress curves of an archive-building search, not as convergence curves from one fixed source genotype.

The current implementation is also greedy at the level of accepted substitutions. A candidate replacement is retained only when it improves the current genotype according to the evaluated objective. This makes the search behaviour simple and easy to interpret, but it can also limit exploration. Potentially useful neutral or temporarily harmful substitutions are discarded, even if they could enable later improvements. As a result, CoSO may converge locally and repeatedly restart rather than traverse low-fitness intermediate regions.

Finally, the experimental results show that building a large archive is not sufficient by itself. Some configurations produced many archive entries or convergence events without producing strong final improvements. The usefulness of the archive depends on whether retrieved fragments are compatible in a way that leads to better complete genotypes. CoSO is therefore best understood as a promising, interpretable genotype-space mechanism whose effectiveness depends on the quality of the fragment descriptors, the chosen compatibility thresholds, and the structure of the available genotype pool.

% TODO: insert discussion after CoSO details/results are available. -DONE - nwm czy nie za mało szczegółowe

\section{Strengths and limitations of Latent Space Optimization}

The principal strength of Latent Space Optimization is the separation between representation learning and search. A trained autoencoder converts a variable-length \fone{} genotype into a fixed-dimensional vector, after which different derivative-free optimizers can be applied without changing the final objective. The latent-space EA, CEM, and CMA-ES all used the same decode--validate--evaluate interface and received true Framsticks \vertpos{} feedback for generated candidates, although their handling of the stored source differed as discussed below. This decoupling made it possible to compare several search strategies on the same learned representation and, conversely, to reuse one search strategy across several checkpoint families.

Grammar-based representation is another important strength. Production sequences, derivation context, and grammar masks provide explicit structure that is unavailable to an unconstrained character decoder. They make syntactic admissibility testable and can restrict locally incompatible production choices before simulator evaluation. The benefit is nevertheless conditional rather than absolute. In the unconditional checkpoint diagnostic, generated validity ranged from 29\% to 98\%, even though all ten selected checkpoints represented \fone{} genotypes through grammar-related models. Grammar awareness therefore reduced the representation problem, but did not guarantee a complete admissible derivation, a simulatable organism, or high fitness at every latent point.

Under the uniform source-retention decision policy, the largest within-LSO means came from latent EA combined with TreeVAE-80 or TransformerVAE-90. Their mean retained scores were 1.772 and 1.774, the two largest means among the fifteen evaluated checkpoint--optimizer configurations. This result identifies effective fixed representation--search combinations within the investigated LSO design space; it is not generated-candidate-only performance and does not 
% establish an advantage over the Framsticks EA baseline. The completed LSO runs and standard baseline runs used different simulation chains, so their aligned starting genotypes do not support an ordinary controlled comparison of final candidate scores. The evidence for LSO is therefore configuration-dependent and must be interpreted without a same-condition baseline claim.
% -- poprawka F
by itself establish an advantage over the Framsticks EA baseline without the paired statistics defined earlier.

This result identifies effective fixed representation--search combinations within the investigated LSO design space; it is not generated-candidate-only performance and, by itself, is not a baseline comparison. The completed LSO runs and the standard Framsticks EA baseline use the same four-file simulation chain and the same starting-genotype manifest, so LSO can be compared with this standard baseline at the paired source-genotype level. This does not apply to the mini baseline or to CoSO: the mini baseline uses a different first simulation file, while CoSO has an archive-based multi-start experimental unit, so direct comparison with LSO will be conducted in different way.

A central limitation is checkpoint dependence. Across the five checkpoints, the latent-space EA retained mean ranged from 1.439 to 1.774. CEM and CMA-ES also changed substantially with the representation. CEM had the lowest optimizer-wide retained mean and was particularly weak with TreeVAE-80 and TransformerVAE-90, although it was more competitive for some LHS checkpoints. CMA-ES reached high maxima, including the largest LSO result of 2.353, and had the largest retained mean for the LHS-depth and fitness-aware checkpoints; these outcomes were nevertheless accompanied by lower central performance than the two leading latent-EA configurations or substantial dispersion in some configurations. Selecting an optimizer without jointly validating the checkpoint is therefore unlikely to be reliable.

The reconstruction results explain why checkpoint selection cannot rely on autoencoder accuracy alone. The LHS latent-256 checkpoint reconstructed 30.81\% of its validation examples exactly, whereas TransformerVAE-90 reconstructed only 2.20\%. Nevertheless, the latter had a markedly larger retained mean with latent EA. Exact reconstruction rewards faithful inversion near observed examples, while optimization also depends on neighbourhood smoothness, the diversity of decoded perturbations, and whether movement in latent space reaches useful phenotypes. The experiments measure these properties only indirectly, but they clearly show that good reconstruction is not sufficient for a strong retained optimization outcome.

The completed protocol also has a source-retention asymmetry. Latent EA inserts the unchanged genotype with its stored manifest fitness and can return that stored-label individual as its raw best; this occurred in 83 of 750 runs. CEM and CMA-ES instead return freshly evaluated decoded candidates and have no equivalent raw fallback. Applying $\max(F^{\mathrm{raw}},F^{\mathrm{source}})$ creates a uniform decision policy and changes the leading optimizer for the LHS-depth and fitness-aware checkpoints, but it cannot reconstruct generated-candidate-only values for the latent-EA runs in which the source remained best. The retained metric is therefore appropriate for the practical decision to keep or replace the source, while raw cross-optimizer rankings from the completed protocol are not methodologically uniform. A future experiment should evaluate the source through the same chain for every optimizer and record generated-candidate-only and retained outcomes separately.

Validity and duplication create an additional search cost. The unconditional generation results ranged from 55\% to 100\% unique strings, and high reconstruction did not imply high random-generation validity. The latent-space EA outputs also recorded large checkpoint-dependent retry counters. Across 300-generation runs, the median invalid-candidate count ranged from 2,137 for TransformerVAE-90 to 92,624 for the latent-256 LHS checkpoint; median duplicate counts ranged from 9 to 19,390. These counters are available for the latent-space EA only and should not be treated as directly comparable CEM or CMA-ES measurements. They nevertheless show that some decoders required many proposals to obtain the requested accepted offspring, reducing the practical efficiency of the learned representation.

LSO also retains the main computational cost of black-box evolutionary design. Neural encoding and decoding do not replace phenotype evaluation: candidate fitness is still obtained from Framsticks. With populations of 100 and 300 generations or iterations, a single run has approximately 30,000 nominal candidate proposals, or evaluation slots. This nominal budget is not the number of unique simulator evaluations: invalid candidates, repeated genotypes, retries, and cached values make both the proposal count and the number of unique evaluated genotypes method-dependent. The autoencoder adds training, checkpoint storage, and inference overhead on top of this simulator cost. Its value lies in proposing a useful search geometry, not in removing the expensive objective evaluation.

Finally, the evidence supports only limited generalization. The final LSO benchmark used five checkpoints, three optimizers, one \fone{} dataset, one fitness criterion and 150 structured starting genotypes. Results may change for another genotype language, simulator configuration, fitness objective, training distribution, or latent dimensionality. The Framsticks EA baseline uses the four-file chain beginning with \texttt{eval-allcriteria-mini.sim}. Broader conclusions require repeated training seeds, independently selected checkpoints, additional tasks, and evaluation conditions that remain identical across the methods being compared.


\section{Comparative discussion}

CoSO and LSO address the disruption caused by uninformed variation in structured genotypes, but they place representation knowledge at different stages of the search. CoSO remains in the discrete \fone{} representation and proposes explicit substitutions of tree-based fragments. Compatibility is determined through hand-designed spatial and directional descriptors, and a substitution is retained greedily only after evaluation of the complete recipient genotype. LSO instead learns a continuous search geometry from genotype data. Its optimizers manipulate latent vectors, while decoding, validity checks, and Framsticks evaluation determine whether a proposed movement becomes a useful \fone{} candidate. The former approach makes individual search moves directly interpretable; the latter can capture regularities that were not explicitly specified, but transfers part of the method design to training, checkpoint selection, and decoding.

The results show that neither representation mechanism should be considered independently of its search policy. CoSO performance changed markedly with the spatial compatibility threshold, indicating that archive utility depends on whether the descriptors retrieve fragments that work in the recipient context. Within LSO, the outcome depended jointly on checkpoint and optimizer. TreeVAE-80 and TransformerVAE-90 produced the largest retained means when combined with latent EA, whereas CMA-ES led for two other checkpoints. Moreover, reconstruction accuracy did not predict this ranking. These findings support a common interpretation: the outcomes associated with representation-aware candidate generation depend on how the representation is traversed and how proposed candidates are filtered and accepted.

Both approaches produced improvements relative to their own references. CoSO improved beyond the best unmodified genotype introduced during an archive-based run, with the strongest aggregate outcomes obtained for spatial thresholds 0 and 1. Within LSO, latent EA with TreeVAE at epoch 80 and with TransformerVAE at epoch 90 produced the largest retained outcomes. At bucket level, these LSO configurations had larger mean improvement margins than the selected CoSO configurations in the completed experiments. This numerical pattern is descriptive rather than evidence that LSO simply outperforms CoSO: an LSO margin is measured from one stored source label under a keep-or-replace policy, whereas a CoSO margin is measured from the best unmodified genotype encountered in a multi-start run. The experimental units, score references, and trajectories are consequently different.

The mini Framsticks EA provides additional context but does not resolve this comparability limitation. CoSO and mini EA evaluate generated candidates with the same four-file mini simulation chain, yet CoSO searches a bucket-level pool with an accumulating archive while mini EA performs one independent run per manifest source. Their standalone aggregates therefore do not constitute an ordinary paired comparison. LSO uses the same source manifest as mini EA, but its generated candidates were evaluated with \texttt{eval-allcriteria.sim} only rather than with \texttt{eval-allcriteria-mini.sim}, \texttt{deterministic.sim}, \texttt{sample-period-2.sim}, and \texttt{only-body.sim}. The LSO--mini plots consequently provide cross-configuration descriptive context and cannot establish superiority, inferiority, or equivalence.

The validity evidence likewise reflects different failure modes. CoSO's tree-based fragments and explicit compatibility criteria reduce arbitrary structural disruption, but direct operation in genotype space does not ensure that every substitution is acceptable: candidates may still require repair or fail evaluation. LSO can constrain production choices through grammar-aware decoding, but validity depends on the checkpoint and on the region visited by the optimizer. The unconditional and latent-EA diagnostics demonstrate substantial variation in valid output, uniqueness, and retry cost. Because comparable optimization-stage counters were not recorded for every method, the evidence is strongest for diagnosing LSO representations rather than for ranking CoSO and LSO by one common validity rate.

The comparative outcome is therefore a trade-off rather than a universal ordering. CoSO offers auditable fragment-level interventions without a learned model, at the cost of domain-specific descriptors and greedy local search. LSO provides reusable continuous optimizers and can learn a search geometry associated with strong outcomes in selected configurations, but its effectiveness depends on the training distribution, checkpoint, optimizer, decoding validity, and simulator-facing inference pipeline. In the studied setting, both mechanisms showed larger measured improvement margins for weak or moderate initial material and smaller margins for the strongest starting range. This is a consistent qualitative within-method pattern in the completed results, although its magnitude cannot be compared as a controlled method effect under the completed protocols.

\section{Threats to validity}

Several aspects of the measurement protocol constrain the interpretation of the results. First, the retained LSO endpoint combines fresh search results with stored source labels. Latent EA included the original genotype with its manifest fitness and returned it as the best individual in 83 of 750 runs. For those source-winning rows, the best score produced only by generated latent-EA candidates is unavailable. CEM and CMA-ES returned freshly evaluated decoded candidates and received source retention only during analysis. The uniform maximum of the raw result and source label is meaningful as a keep-or-replace decision, but it is not a uniform generated-candidate-only optimizer measurement.

The stored source labels create a related construct-validity limitation. They determine bucket membership and the reference for retained improvement, but generated candidates receive fresh simulation values. A stored label and a fresh result should not automatically be interpreted as repeated measurements under identical conditions. This concern is particularly important for the LSO--mini contrast: completed LSO candidates used \texttt{eval-allcriteria.sim} only, whereas mini EA used \texttt{eval-allcriteria-mini.sim} followed by \texttt{deterministic.sim}, \texttt{sample-period-2.sim}, and \texttt{only-body.sim}. Sharing the same manifest aligns source identities, not the score-producing conditions. Accordingly, the reported LSO--mini differences remain a cross-configuration descriptive contrast rather than controlled paired effects.

Comparisons involving CoSO have a separate unit-of-analysis limitation. CoSO and mini EA share the mini simulation chain, but CoSO is an archive-based multi-start process over a genotype pool, whereas mini EA is replicated once for each individual source. Similarly, CoSO improvement is referenced to the best unmodified genotype introduced in a run, while LSO improvement is referenced to one stored source genotype. Aggregate and bucket-level patterns can be contrasted, but treating those observations as paired replicates would overstate the design. A large CoSO archive or a high number of local convergence events also measures search activity rather than solution quality by itself.

The nominal budgets do not imply equal computational effort. LSO and mini-EA runs used approximately 30,000 requested candidate slots per source, whereas CoSO used a limit of 30,000 within-constraint evaluations per run. Invalid outputs, duplicate genotypes, retries, cache hits, initialization, and repair affect the relation between these budget axes and the number of distinct simulator calls. The available counters are method-dependent: detailed invalid and duplicate proposal counts are available for latent EA, but not on a directly comparable basis for CEM, CMA-ES, CoSO, and mini EA. Runtime and simulator-call comparisons are therefore not supported by a common measurement definition.

Model development introduces further internal and conclusion-validity risks. Autoencoder validation results came from run-specific 20\% splits; the split size was common, but membership cannot be assumed to be identical across all training runs. The final checkpoint set was selected from a broader development process, and the final optimization used fixed checkpoints and deterministic seed derivation rather than repeated independent training and optimization campaigns for every configuration. Twenty-five runs per bucket characterize variation over the selected sources and run seeds, but they do not quantify uncertainty due to retraining a model, choosing another held-out split, or repeating checkpoint selection. Consequently, observed checkpoint--optimizer rankings may include selection-specific effects.

The statistical conclusions are intentionally descriptive. Means, medians, dispersions, maxima, retention counts, and trajectories summarize the completed runs, but no inferential procedure establishes population-level differences between methods. This restraint is necessary because score references, experimental units, simulation chains, and available diagnostic counters are not uniform across all conditions. High maxima, in particular, should not be interpreted as evidence of stable performance, and larger selected-configuration bucket margins should not be generalized into method-wide dominance.

External validity is bounded by the empirical scope. The experiments use one dataset of \fone{} genotypes, one phenotype simulator, one objective (\vertpos{}), one set of bucket boundaries, and a limited collection of CoSO descriptors, autoencoder checkpoints, and optimizer settings. Other objectives may reward different structures, and other genotype representations may not admit the same fragment descriptors or grammar. The results therefore establish behaviour for the investigated Framsticks design task rather than a general ranking of explicit substitution and learned latent-space optimization.

\section{Conclusions}

This thesis investigated two representation-aware approaches to evolutionary design of variable-length \fone{} genotypes. The empirical evidence answers the active research questions within the limits of the completed protocols.

For \textbf{RQ1}, the experiments establish that both approaches can produce evaluable improved genotypes, but answer comparative reliability only partially. CoSO substitutes coherent fragments selected through spatial and directional compatibility, while still relying on repair and complete-candidate evaluation. In LSO, grammar-aware decoding makes admissibility explicit, yet unconditional validity varied substantially between checkpoints and latent EA incurred strongly checkpoint-dependent invalid and duplicate retry counts. Because equivalent optimization-stage validity and evaluability counters are unavailable for all methods, a common reliability rate cannot be determined. The most detailed conclusion therefore concerns LSO: decoder validity and diversity must be evaluated separately from reconstruction accuracy and final fitness.

For \textbf{RQ2}, the compared methods improved starting material according to their respective reported references. The TreeVAE-80 and TransformerVAE-90 latent-EA configurations had mean retained improvements of approximately 0.763 and 0.764 over their stored source labels. CoSO spatial thresholds 0 and 1 had mean run-level improvements of 0.465 and 0.455 over the best unmodified genotype introduced into each archive-based run. The mini Framsticks EA had a standalone mean source-reference improvement of 0.683. These magnitudes do not form a common-condition ranking: CoSO uses a different experimental unit and reference, while LSO and mini EA use different simulation configurations. Moreover, the retained LSO endpoint is a keep-or-replace outcome and does not recover generated-only performance for all latent-EA runs.

For \textbf{RQ3}, explicit compatible substitution and learned latent-space search exhibit complementary advantages. CoSO provides interpretable fragment replacements and transparent compatibility control without training a generative model. For the selected leading configurations, LSO produced larger bucket-conditioned improvement margins in the completed experiments. That observation is not a direct dominance result because CoSO and LSO use different experimental units and reference scores. The defensible comparison is that CoSO offers greater move-level interpretability and lower representation-training overhead, whereas LSO provides a learned continuous search space whose observed outcomes are more sensitive to checkpoint, optimizer, decoding, and inference choices.

For \textbf{RQ4}, initial fitness was strongly associated with the measurable room for improvement. CoSO and every fixed LSO configuration had larger mean improvement in the weakest bucket than in the strongest bucket, although intermediate LSO buckets were not always monotonic. Strong starting genotypes more often triggered source retention in LSO and yielded smaller CoSO gains over the best introduced genotype. Under the reported metrics, weak or medium-quality material was therefore associated with larger improvement margins than already strong designs.

The bounded thesis-level conclusion is that representation-aware search is useful for structured evolutionary design, but the representation alone does not determine success. CoSO depends on descriptors, archive retrieval, and acceptance policy; LSO depends on learned geometry, checkpoint--optimizer compatibility, decoding validity, and true simulator evaluation. The completed evidence supports these within-method findings and an aggregate comparison of their tendencies. It does not establish a controlled ranking against the mini Framsticks EA or a universal ordering between CoSO and LSO.

\section{Future work}

The first priority is a controlled comparison protocol. CoSO, LSO, and mini Framsticks EA should evaluate every generated candidate with an identical complete simulation chain and should apply the same treatment to source genotypes. Each run should record both the best generated-only result and the final result after an explicit source-retention decision, with the source freshly evaluated under that chain. For CoSO, a single-source variant would provide the clearest pairing with LSO and mini EA. If archive-based multi-start remains essential, the other methods should instead be evaluated on matched pools with an equivalent run-level reference and independently reset state. This design would directly address the present simulation-chain, source-semantics, and experimental-unit limitations.

A second priority is to quantify uncertainty introduced by representation learning and configuration selection. Autoencoders should be retrained with several independent seeds and declared held-out splits, followed by repeated checkpoint selection and optimization. CoSO threshold selection and optimizer hyperparameter selection should likewise be separated from final assessment. Reporting results on a final untouched test manifest would distinguish sensitivity to sources from sensitivity to training, selection, and optimization randomness.

Method development for CoSO should target the limitations of greedy, hand-designed compatibility. Neutral or temporarily detrimental substitutions could be admitted through probabilistic, annealed, or population-based acceptance to test whether crossing local fitness valleys improves final outcomes. Spatial and directional descriptors could be augmented with richer structural, behavioural, or learned fragment features, while retaining an interpretable account of why a donor fragment matches its recipient context. Ablations of repair, archive composition, descriptor components, and retrieval thresholds would show which parts of the mechanism contribute to valid and useful substitutions.

For LSO, future models should optimize not only reconstruction but also the properties required by search. Suitable diagnostics include local decode stability, validity under optimizer-induced perturbations, neighbourhood diversity, phenotype continuity, and the relation between latent distance and changes in \vertpos{}. Training objectives or decoding procedures could then explicitly discourage invalid and duplicate regions. These experiments should preserve true Framsticks evaluation as the final objective and should report generated-only and retained results separately for latent EA, CEM, and CMA-ES.

The complementary strengths of the methods motivate a hybrid design. CoSO fragments could define interpretable edits whose donor selection or compatibility metric is informed by a learned embedding. Conversely, latent search could propose complete genotypes that enrich a CoSO archive, after which explicit substitutions refine or recombine their useful substructures. A controlled ablation should compare each standalone component with the hybrid under matched source handling and simulator evaluations, so that any gain can be attributed to interaction rather than to a larger effective budget.

Finally, the empirical scope should be broadened in response to the external-validity and measurement limitations. Further studies should include other Framsticks objectives, multi-objective criteria, additional genotype representations, and out-of-distribution source sets. Solution quality should be complemented by structural and behavioural diversity, genotype and phenotype complexity, robustness to simulation perturbations, wall-clock time, repair frequency, and exact counts of unique simulator evaluations. Standardized manifests, simulation-chain identifiers, generated/retained endpoints, candidate-status counters, and machine-readable run histories would make these comparisons reproducible and prevent nominal candidate slots from being mistaken for equivalent computational effort.
